\section{Нахождение обратной матрицы методом элементарных преобразований}

\begin{theorem}
Квадратная матрица $A$ порядка $n$ путем элементарных преобразований только строк может быть приведена к единичной матрице $E$ тогда и только тогда, когда $\det A \neq 0$.
\end{theorem}

\begin{proof}
\begin{itemize}
    \item \textbf{Необходимость ($\Rightarrow$):} Пусть $A$ элементарными преобразованиями строк приводится к $E$. Это означает, что существует последовательность матриц элементарных преобразований $B_1, B_2, \dots, B_k$ такая, что $B_k \dots B_2 B_1 A = E$. Перейдя к определителям, получим:
    $$
    \det(B_k) \dots \det(B_2) \det(B_1) \det(A) = \det(E) = 1.
    $$
    Так как матрицы элементарных преобразований невырожденны ($\det B_i \neq 0$), то и $\det A \neq 0$.
    \item \textbf{Достаточность ($\Leftarrow$):} Пусть $\det A \neq 0$. Тогда $\operatorname{rang} A = n$. По теореме о ступенчатом виде, $A$ можно привести к ступенчатой матрице. Поскольку ранг равен $n$, в ступенчатой матрице нет нулевых строк, и все ведущие элементы находятся на главной диагонали и отличны от нуля. Умножая строки на обратные величины и прибавляя их к вышележащим строкам, мы можем обнулить все элементы выше главной диагонали и сделать диагональные элементы равными $1$. Таким образом, $A$ приводится к $E$.
\end{itemize}
\hfill$\blacksquare$
\end{proof}

\begin{theorem}
Если квадратная матрица $A$ невырождена ($\det A \neq 0$), то, применяя к расширенной матрице $(A \mid E)$ одни и те же элементарные преобразования строк, мы приведем ее к виду $(E \mid A^{-1})$.
\end{theorem}

\begin{proof}
Приведение $A$ к $E$ означает, что существует произведение матриц элементарных преобразований $B = B_k \dots B_1$ такое, что $BA = E$. Отсюда следует, что $B = A^{-1}$.
Если мы применим те же самые преобразования (то есть умножим слева на $B$) к расширенной блочной матрице $(A \mid E)$, то в силу правил умножения блочных матриц получим:
$$
B \cdot \left( \begin{array}{c|c} A & E \end{array} \right) = \left( \begin{array}{c|c} BA & BE \end{array} \right) = \left( \begin{array}{c|c} E & A^{-1} \end{array} \right).
$$
Следовательно, в правой части расширенной матрицы автоматически окажется обратная матрица $A^{-1}$.
\hfill$\blacksquare$
\end{proof}

\begin{example}
Для нахождения $A^{-1}$ записывают блочную матрицу и приводят левый блок к единичному:
$$
\left( \begin{array}{c|c}
A & E
\end{array} \right) 
\xrightarrow{\text{эл. преобр. строк}} 
\left( \begin{array}{c|c}
E & A^{-1}
\end{array} \right)
$$
\end{example}
